Recordemos que $A^*$ es un algoritmo de búsqueda informada definido por
\[
f(n)=g(n)+h(n),
\]
donde $g(n)$ es el costo acumulado hasta $n$ y $h(n)$ es la heurística asociada.

Por otro lado, un algoritmo de búsqueda es completo si encuentra una solución siempre que esta exista.\\

\boxed{\textbf{Dem.}}

Supongamos que existe al menos un nodo meta alcanzable desde el estado inicial, que el índice de ramificación del árbol de búsqueda es finito ($b<\infty$) y que todos los costos de acción son positivos.

Sea $C^*$ el costo de una solución óptima.

Sabemos que $A^*$ siempre elige para expandir el nodo $n$ de la frontera con menor valor de
\[
f(n)=g(n)+h(n).
\]

Entonces, cualquier nodo que cumpla
\[
f(n)<C^*
\]
va a ser expandido antes que cualquier nodo con valor
\[
f(n)\geq C^*,
\]
simplemente porque $A^*$ siempre prioriza el menor $f$ disponible en la frontera.\\

Ahora bien, como el índice de ramificación es finito y los costos son positivos, al bajar en el árbol el costo acumulado $g(n)$ va creciendo. Eso implica que no puede haber infinitos nodos con costo estimado estrictamente menor que $C^*$; en otras palabras, no pueden existir infinitos nodos tales que
\[
f(n)<C^*.
\]

Por lo tanto, el conjunto de nodos que satisfacen
\[
f(n)<C^*
\]
es finito.\\

Como $A^*$ expande nodos de menor a mayor $f(n)$, primero va a sacar de la frontera todos los nodos que cumplan
\[
f(n)<C^*.
\]

Y como ya vimos que esos nodos son finitos, ese proceso termina.

Entonces, después de expandirlos a todos, solo pueden quedar en la frontera nodos con
\[
f(n)\geq C^*.
\]

Como existe una solución óptima de costo $C^*$, en algún momento el nodo meta óptimo será generado y quedará en la frontera.

Además, como $A^*$ siempre expande el nodo con menor valor de $f(n)$, ese nodo meta solo puede ser pospuesto por nodos con
\[
f(n)<C^*.
\]

Pero ya vimos que la cantidad de nodos con
\[
f(n)<C^*
\]
es finita.

Por lo tanto, después de expandir todos esos nodos, $A^*$ terminará seleccionando el nodo meta.

Así, si existe una solución, $A^*$ eventualmente encuentra una meta.

\[
\therefore\ A^* \text{ es completo.}
\]
